% !TEX root = ../../main.tex
% Part II (Nguyen Huu Cuong) - merged from the research report (report_latex_v3), see MERGE_NOTES.md.
% Style pass 2026-09-30: em dashes and rhetorical sentences removed; technical content, numbers and references unchanged.
% Tightened 2026-10-05 (both authors agreed): shorter sentences; the "why not CLIP alone" argument is stated once (9.1.2).
\chapter{Semantic Reasoning: Motivation and Overview}
\label{chapter:r-overview}

\textit{This chapter opens Part~II, the research on the search block. It states the problem and the contributions, reads the person search pipeline as a template of replaceable blocks, argues for improving the search block with a reasoning layer inspired by ZARA, and fills each block with a concrete model.}

\section{Motivation and Contributions}
\label{sec:r-intro}

\subsection{Problem}
Given a free-text description such as \emph{``a teenage woman, with short hair,
wearing a short-sleeve red top and shorts, with a bag''} and a gallery of person
crops from surveillance video, the task is to return the crops that show that
person. This is text-based person search (Section~\ref{sec:3.1}), here over small,
blurred, badly lit crops taken from arbitrary viewpoints.

\subsection{Why a reasoning layer}
A dual-encoder model such as CLIP encodes the sentence and every crop and sorts by
cosine similarity. This is fast and the right first stage, but it has two limits.

\begin{itemize}
\item It gives one opaque number per crop. When the right person is ranked 37th,
      the score cannot say whether the top colour, the sleeves or the gender was
      misread, so nothing can be corrected, and fine-tuning only changes the number.
\item It ignores the track. Tracking groups crops of the same person seen under
      different blur, lighting and pose, but CLIP scores each crop alone and
      discards that redundancy.
\end{itemize}

An attribute-level representation addresses both: it is readable and can be
aggregated across the frames of a track. The attributes, however, come from a small
VLM and are noisy (Section~\ref{sec:r-noise}), so they are aggregated by an LLM
that reasons over evidence rather than by a majority vote.

\subsection{Contributions}
\begin{enumerate}
\item A statement of where the improvement should go: the pipeline is read as a
      template of replaceable blocks, and the search block is chosen because it
      must both generalise semantically and tolerate the errors of every block
      upstream (Section~\ref{sec:r-directions}).
\item An attribute-centric index over tracked person crops: 18 attributes per crop
      read by Qwen2.5-VL-3B, with image-quality statistics per body region
      (sharpness, contrast, brightness, saturation) that let a later stage reason
      about \emph{why} two readings disagree.
\item An adaptation of the ZARA evidence-and-reasoning pattern from motion-sensor
      activity recognition to image-text person retrieval
      (Section~\ref{sec:r-zara}).
\item A four-stage multi-agent funnel (query analyst, feature selector, evidence
      reader, ranker) over two rounds with different cut semantics, with the prompt
      mechanisms that make it work: a citation rule, a reliability-not-direction
      ranking rule, a rarity table and a round-dependent coverage floor.
\item An evaluation protocol and tooling that rebuild every ranking from the saved
      model replies, and that separate retrieval from reasoning failures by
      reporting the pool ceiling next to every recall number.
\item A measured result on 100 identities, decomposed into what the layer rescues
      and what it costs, with an error analysis of both.
\end{enumerate}

% ============================================================= 2 background =
\section{The Pipeline, and Where to Improve It}

\subsection{The pipeline in the abstract}
\label{sec:r-pipeline}
Figure~\ref{fig:r-overall} names every block by what it \emph{does} rather than by
what implements it: a person search system has five replaceable parts and two
phases. The question is then which block a fixed budget should go into.

\begin{figure}[H]
\centering
\includegraphics[width=\textwidth]{S23-overall-pipeline.png}
\caption[The pipeline as a template of replaceable blocks]{The pipeline as a template. Offline: video $\rightarrow$ object
detection $\rightarrow$ tracker $\rightarrow$ image encoder $\rightarrow$
database. Online: text query $\rightarrow$ text encoder $\rightarrow$ similarity
search against the same database. Every box is a slot for any detector, tracker
or encoder.}
\label{fig:r-overall}
\end{figure}

\subsection{Possible directions for improvement}
\label{sec:r-directions}
Table~\ref{tab:r-directions} compares the blocks by what each must achieve and how
it fails.

\begin{table}[H]
\centering
\begin{tabular}{@{}P{3.2cm}P{4.0cm}P{6.4cm}@{}}
\toprule
\textbf{Direction} & \textbf{Main requirement} & \textbf{Our consideration}\\
\midrule
Object detection & Detect people correctly & Errors propagate downstream.\\
\addlinespace
Tracking & Maintain identity across frames & Tracker errors can contaminate a
whole track.\\
\addlinespace
Image/text encoders and online search & Match the query to visual appearances &
Needs semantic generalisation and robustness to upstream errors, at the image
quality of deployed cameras.\\
\bottomrule
\end{tabular}
\caption[The three places where the pipeline can be improved]{The three places the pipeline can be improved, what each has to do, and
what that costs.}
\label{tab:r-directions}
\end{table}

Detection and tracking are well studied, and improving them mostly means a better
model in the slot. Their failures also reach the search block anyway: a missed
detection is a crop that never exists, and a merged track contains two people.

The search block alone must generalise from a sentence to an appearance, tolerate
the detector's and tracker's mistakes, and do both on deployed-camera crops; it is
also where the answer is decided. The work of this part is therefore placed there.
Given the two limits of CLIP above, the block needs a representation that can be
\emph{read}, attributes with their evidence, and a component able to weigh several
partly contradictory readings of one person. This is a reasoning problem over text,
which suits an LLM, as ZARA observed in another domain.

\subsection{ZARA: related work and inspiration}
\label{sec:r-zara}
The closest prior work is ZARA~\cite{li2025zara}, from which the reasoning pattern
is taken; the domain and the evidence are this project's own. ZARA's key idea is
that an LLM reasons poorly over raw numerical signals but well over \emph{textual}
evidence, so it first converts statistical patterns of the reference data into a
textual knowledge base and, at inference, retrieves the discriminative evidence for
the candidate classes (Figure~\ref{fig:r-zara}).

We keep this structure and change the domain. ZARA's statistics come from motion
signals; ours are human-readable person attributes extracted offline by Qwen-VL.
In ZARA, several sensors give complementary views of one activity; here, the
frames of one track give complementary views of one person. The difference is that
a tracker can group two people in one track, an upstream failure that ZARA's
sensors do not have (Section~\ref{sec:r-directions}).

The attribute statistics are not treated as ground truth, since Qwen-VL errs,
especially on real surveillance image quality and viewpoints, and the tracker adds
its own errors. They are given to an LLM as \emph{evidence to be reasoned over},
and the aggregated track-level reading re-ranks the candidates against the query:

\begin{center}
\small
\begin{tabular}{@{}ll@{}}
\toprule
ZARA & motion signals $\rightarrow$ statistical features $\rightarrow$
       multi-sensor evidence $\rightarrow$ LLM reasoning\\
Ours & person images $\rightarrow$ human-readable attributes $\rightarrow$
       multi-frame track evidence\\
     & \quad $\rightarrow$ LLM evidence aggregation $\rightarrow$ candidate
       re-ranking\\
\bottomrule
\end{tabular}
\end{center}

\begin{figure}[H]
\centering
\includegraphics[width=\textwidth]{S26-zara.png}
\caption[ZARA's retrieval of discriminative evidence]{ZARA's retrieval of discriminative evidence for candidate classes, the
pattern this project adapts.}
\label{fig:r-zara}
\end{figure}

The coarse-to-fine iteration is also taken from ZARA: the funnel of
Section~\ref{sec:r-reasoning} prunes before it decides, with a coarse round over the
whole pool and a finer round over the survivors.

\section{System Overview}

\subsection{Filling the slots}
\label{sec:r-choices}
Table~\ref{tab:r-choices} puts a model into each block of Figure~\ref{fig:r-overall}.
The detector and tracker are established defaults, chosen so that they are not the
bottleneck, and are not studied further. In the search block, CLIP stays as the
\emph{base encoder} that builds the candidate pool, followed by an LLM-based
reasoning layer.

\begin{table}[H]
\centering
\begin{tabular}{@{}lll@{}}
\toprule
\textbf{Block} & \textbf{Choice} & \\
\midrule
Object detection            & YOLO11n                     & \cite{jocher2024yolo11}\\
Tracker                     & ByteTrack                   & \cite{zhang2022bytetrack}\\
Image/text encoder          & CLIP ViT-H/14 (laion2B), as the base encoder
                                                          & \cite{laion2022clipvith14}\\
Online search               & LLM-based reasoning layer over the CLIP pool & \\
Database, similarity search & Milvus                      & \cite{milvusdocs}\\
VLM (attribute describer)   & Qwen2.5-VL-3B-Instruct      & \cite{qwen2024qwen25vl3b}\\
LLM (reasoning)             & Gemini Flash-Lite 3.5       & \cite{googledeepmind-geminiflashlite}\\
\bottomrule
\end{tabular}
\caption{One model per slot. None of them is fine-tuned on this dataset.}
\label{tab:r-choices}
\end{table}

\subsection{Offline indexing}
\label{sec:r-offline}
In the indexing phase (Figure~\ref{fig:r-offline-new}), detection and tracking are
unchanged; the new element is that each tracked crop passes through \emph{two}
describers: CLIP for the retrieval embedding and Qwen2.5-VL-3B for the attribute
record. Both are stored under (video, track, frame), so a retrieved crop is joined
with its attributes and with the other frames of its track in one lookup.

\begin{figure}[H]
\centering
\includegraphics[width=\textwidth]{S31-offline-indexing.png}
\caption[Extended offline indexing with the attribute describer]{Extended offline indexing. The Qwen2.5-VL-3B branch turns each crop into
18 readable attributes alongside its embedding, which the reasoning stages need.}
\label{fig:r-offline-new}
\end{figure}

\subsection{Online search}
\label{sec:r-online}
In the search path (Figure~\ref{fig:r-online-full}), CLIP retrieval is unchanged,
but its output becomes an \emph{intermediate candidate pool}, which the multi-agent
reasoning block takes together with the original text to return the final top-$K$.

\begin{figure}[H]
\centering
\includegraphics[width=0.95\textwidth]{S35-text-retrieval.png}
\caption[Online search with the reasoning layer]{Online search with the reasoning layer. CLIP produces the candidates and
the funnel orders them. The user's text reaches the reasoning block directly and
is not summarised by the retrieval stage.}
\label{fig:r-online-full}
\end{figure}

The reasoning layer can only reorder what retrieval gave it: a person missing from
the pool cannot be returned. Section~\ref{sec:r-metrics} therefore reports this
ceiling, \code{cap@75}, next to every recall number.

% ============================================================ 4 the index ===
